\documentclass[11pt]{article}
\usepackage{amsmath,amssymb,amsthm}
\newtheorem{lemma}{Lemma}
\newtheorem{proposition}{Proposition}
\title{Balancing a Shared Cauchy Codebook}
\author{}
\date{November 23, 2026}
\begin{document}
\maketitle

\paragraph{Status and scope.}
This note removes a multiplicity issue in a proposed joint Cauchy decoder.
It does not remove the decoder's cross-block interference, prove a
superconstant local saving, or prove either requested asymptotic bound.
The statements below are complete written proofs and have not been
formalized in Lean. No model definition is changed.

\begin{lemma}[A common relabeling balances every target column]
Let $B,D,s\geq1$ be integers and let
$a_{bj}\in\mathbb Z/D\mathbb Z$ for $1\leq b\leq B$, $1\leq j\leq s$.
There exist shifts $\rho_1,\ldots,\rho_B\in\mathbb Z/D\mathbb Z$ such that
for every column $j$ and every color $d$,
\[
 \#\{b:a_{bj}+\rho_b=d\}\leq
 L:=\left\lceil\frac{2B}{D}+\log(sD)+1\right\rceil.
\]
The shifts are chosen once for the complete table, not separately for
each column.
\end{lemma}
\begin{proof}
Choose the shifts independently and uniformly. For fixed $(j,d)$ the
occupancy $X$ is a sum of $B$ independent Bernoulli variables of mean
$1/D$. Therefore
\[
 \mathbb E e^X=\left(1+\frac{e-1}{D}\right)^B
 \leq \exp\left((e-1)B/D\right)<\exp(2B/D).
\]
Markov's inequality gives
\[
 \Pr(X\geq L)<\exp(2B/D-L)\leq\frac1{e sD}.
\]
The union bound over the $sD$ pairs gives probability less than $1/e$
that any occupancy is at least $L$. Hence some shifts have every
occupancy less than $L$, which is stronger than the claimed inequality.
\end{proof}

\begin{proposition}[Balanced assignments with linear overhead]
In the lemma, put $G=DL$. There are $G$ anchor slots and $G$ color slots,
exactly $L$ of each color. Each of the $B$ original anchors in column $j$
can be assigned its prescribed shifted color $a_{bj}+\rho_b$, and the
remaining $G-B$ anchors can be assigned so that the full assignment is
balanced. If $D(\log(sD)+2)\leq B$, then $G\leq3B$.
\end{proposition}
\begin{proof}
The number of occupied slots of each color is at most $L$. Fill its
remaining slots with unused dummy anchors. The total number of vacancies
is $DL-B=G-B$, so this exhausts precisely the dummy anchors. The ceiling
bound gives
\[
 G=DL\leq 2B+D\log(sD)+2D\leq3B.
\]
Also $G\geq2B$, so no negative number of dummy anchors occurs.
In the grouped Cauchy transfer, give each dummy anchor residue zero.
It contributes zero to every grid query. The existing paired construction
then uses $2G\leq6B$ heads before the established constant-cost axis
repair. This statement concerns assignment multiplicities only.
\end{proof}

\paragraph{What the codebook proposal does establish.}
Fix an invertible $k$ by $k$ exterior Cauchy matrix $C$. For $q$ target
bit columns, an exact block interpolant has coefficient array
$W=C^{-1}S$ with $S\in\{-1,1\}^{k\times q}$. Consequently every row
of every such $W$ belongs to a fixed finite codebook of at most
$k2^{kq}$ vectors in $\mathbb R^q$. This is an exact cardinality bound;
it is independent of the condition number of $C$. It does not bound
the magnitudes or numerical stability of these coefficients. The
balancing lemma permits a shared alphabet to be reused across many
blocks without allocating one copy of the full alphabet to each block.
A block's fixed shift can be undone in that block's query profile.

\paragraph{The missing assembly step.}
The grouped family evaluates a query by
\[
 f_i=\sum_{g=1}^G\frac{c_g v_{i,\ell(g)}}{(t_i-\xi_g)^2}.
\]
Undoing the shift for block $b$ also changes the interpretation of every
other block's colors at a row in block $b$. Those contributions remain
present. Inverting the small matrix $C$ controls only the intended
block's terms and supplies no cancellation of the other terms. The
balancing lemma does not claim such cancellation. Thus it cannot be
composed with the block interpolation statement to assert a universal
decoder. Establishing the required simultaneous cancellation, or a
strict-sign estimate that makes it unnecessary, remains open.

\paragraph{Two restricted readout choices cannot resolve that gap.}
If all query poles are one fixed $t$, every assignment and choice of
residues produces a vector $R\in\mathbb R^D$, with query value $v^\top R$.
Thus no query set of size more than $D$ can be shattered: it is contained
in the class of homogeneous linear halfspaces in $\mathbb R^D$.

If instead every query profile is a nonzero multiple of a single
coordinate vector, group the query points by that coordinate. Suppose
each color has $k=G/D$ assigned simple poles. On one color, the readout
\[
 R_d(t)=\sum_{g:\ell(g)=d}\frac{c_g}{t-\xi_g}
\]
has, after multiplication by its positive denominator, a polynomial
numerator of degree at most $k-1$. Distinct sorted exterior query poles
therefore cannot carry more than $k-1$ sign changes. In particular a
shattered set has at most $k$ queries of each color and at most $G$
queries in total. A zero residue or a coincident query cannot increase
this bound. For squared kernels the common denominator has degree $2k$
and the numerator has degree at most $2k-2$, giving at most $2k-1$
queries per color and at most $2G-D$ overall. Nonzero scalar multiples
of coordinate vectors only predetermine sign flips and do not change
the argument. These restrictions leave mixed color profiles with
variable query poles entirely open.

\paragraph{Finite check.}
The companion exact-arithmetic script verifies the shared relabeling,
balanced dummy filling, zero dummy residues, and a small rational Cauchy
block codebook. These tests check the stated constructions, not the
unproved cross-block cancellation or asymptotic universality.
\end{document}
